\section{Related Work}

Genomic antimicrobial resistance prediction is already an active research direction, which itself shows that the problem is meaningful. The central goal is to infer observed resistance phenotype from genomic evidence quickly and at scale, but this goal remains difficult because genotype does not map perfectly to phenotype. Known resistance genes and mutations provide useful biological clues, yet they do not always explain the final laboratory phenotype completely. At the same time, broader genome-wide representations may improve predictive coverage but often weaken direct biological interpretation.

As a result, the literature has not converged on one universally best solution. Some studies emphasize curated biological knowledge and direct interpretability, some use biologically meaningful features together with machine learning, and others move toward broad genome-wide representations such as SNPs, k-mers, or pan-genome features~\citep{feldgarden2021amrfinderplus,maguire2019identification,ren2022prediction}. This diversity is important for the present report because the main scientific question is not only whether AMR can be predicted, but also how the choice of genotype-to-phenotype representation affects the balance between predictive performance and biological explanation.

Table~\ref{tab:related-work-overview} summarizes representative studies that are most relevant to the present project. The purpose of this table is methodological positioning rather than direct numerical benchmarking, because the studies differ in organism, phenotype target, feature representation, and evaluation protocol.

The first family of studies begins from known AMR genes, known resistance-associated mutations, and curated biological annotation. In this direction, the main advantage is interpretability: every prediction or association can be traced back to a specific determinant or mechanism. The AMRFinderPlus framework of Feldgarden et al.~\citep{feldgarden2021amrfinderplus} is especially important here because it provides the curated determinant catalog that motivates the feature representation used in this report. Shen et al.~\citep{shen2024comparison} is also important because it directly compares genotypic and phenotypic resistance profiles in \textit{Salmonella enterica}, which is close in spirit to the rule-based reference baselines used in the present study. The limitation of this family is that direct inference can fail when resistance depends on broader within-class mechanisms, unknown variants, or incomplete determinant coverage.

The second family uses biologically meaningful features but lets machine learning learn the final mapping from genotype to phenotype. This is the closest family to the present project. Maguire et al.~\citep{maguire2019identification} and Benefo et al.~\citep{benefo2024genomebased} both illustrate the value of combining mechanism-level features with supervised learning, so that the model can capture determinant interactions and weighting effects that a fixed rule cannot express easily. He et al.~\citep{he2025prediction} extends this direction toward broader pan-genome and pan-resistome modeling in \textit{Salmonella}, which is particularly relevant when comparing our biologically constrained settings with the wider \textit{all}-scope representation.

The third family moves toward genome-wide or ML-heavy representations. Ren et al.~\citep{ren2022prediction} uses whole-genome sequencing with SNP-based features, while Wang et al.~\citep{wang2023using} predicts MIC from genome-wide 10-mers in nontyphoidal \textit{Salmonella}. These studies show that useful AMR signal may exist outside curated determinant catalogs, and they often achieve strong predictive performance. However, their broader feature spaces make the biological meaning of individual features harder to interpret directly. For the present report, they are therefore best treated as high-coverage comparison directions rather than as directly comparable determinant-based baselines.

The present study belongs primarily to the biology-informed machine-learning family, but it was intentionally designed to sit between strict rule-based inference and broader genome-wide prediction. It uses curated determinant-level features from AMRFinderPlus and MicroBIGG-E for interpretability, adds rule-based baselines to show how far direct genotype-to-phenotype inference can go, and then compares several genotype-to-phenotype connection scopes to measure how predictive gain changes as biological constraints are relaxed.

This position is useful because it avoids two extremes. It is more biologically explicit than a purely genome-wide black-box model, yet more flexible than a direct deterministic rule system. The resulting contribution of the study is therefore not merely another classifier benchmark. Instead, it is a controlled comparison of how curated determinant evidence, biological scope design, and learned combination interact in genomic AMR prediction.

\begin{table}[!htbp]
\centering
\caption{Representative related studies and their relation to the present project}
\label{tab:related-work-overview}
\scriptsize
\setlength{\tabcolsep}{3pt}
\renewcommand{\arraystretch}{1.1}
\begin{tabular}{@{}p{2.1cm} p{1.8cm} p{2.1cm} p{2.0cm} p{4.8cm}@{}}
\toprule
Study & Target & Feature type & Method family & Relation to our study \\
\midrule
Feldgarden et al.~\citep{feldgarden2021amrfinderplus} & General AMR annotation & Curated determinants & Biology-first resource & Foundation for the determinant space used in our pipeline; not a direct prediction benchmark. \\
Shen et al.~\citep{shen2024comparison} & Salmonella concordance & Genotypic AMR profiles & Rule-based direction & Closest conceptual basis for our direct genotype-to-phenotype baseline. \\
Maguire et al.~\citep{maguire2019identification} & Salmonella AMR prediction & AMR features & Biology-informed ML & Close in spirit to our determinant-based machine-learning setting. \\
Benefo et al.~\citep{benefo2024genomebased} & Salmonella AMR prediction & Genome-based features & Biology-informed ML & Similar biological task, but different representation and dataset context. \\
He et al.~\citep{he2025prediction} & Salmonella phenotype and MIC & Pan-genome / pan-resistome & Broader ML & Closest literature analogue to our wider-scope machine-learning direction. \\
Ren et al.~\citep{ren2022prediction} & Binary AMR prediction & Whole-genome SNPs & Genome-wide ML & Useful for discussing broad genomic signal beyond curated determinants. \\
Wang et al.~\citep{wang2023using} & Salmonella MIC prediction & WGS-derived 10-mers & Genome-wide ML & Useful mainly for the richer MIC target and the interpretability trade-off of genome-wide features. \\
\bottomrule
\end{tabular}
\end{table}
